\documentclass[12pt,a4paper]{article}

% ---------- Кодировка и язык ----------
\usepackage{cmap}
\usepackage[T2A]{fontenc}
\usepackage[utf8]{inputenc}
\usepackage[english,russian]{babel}

% ---------- Математика ----------
\usepackage{amsmath}
\usepackage{amssymb}
\usepackage{mathtools}

% ---------- Таблицы и рисунки ----------
\usepackage{graphicx}
\usepackage{booktabs}
\usepackage{tabularx}
\usepackage{multirow}
\usepackage{makecell}
\usepackage{caption}
\usepackage{subcaption}
\usepackage{float}

% ---------- Оформление ----------
\usepackage[a4paper,margin=2.5cm]{geometry}
\usepackage{setspace}
\usepackage{enumitem}
\usepackage{xcolor}
\usepackage{microtype}
\usepackage[authoryear,round]{natbib}
\usepackage[hidelinks]{hyperref}

\graphicspath{{figures/}}
\onehalfspacing
\setlength{\parindent}{1.25cm}
\setlength{\parskip}{0pt}

% ---------- Данные работы ----------
% Рабочие формулировки уточняются по мере получения результатов.
\newcommand{\researchtitle}{Исследование переноса группового знания в нейронных процессах для cold-start рекомендаций}
\newcommand{\researchtitleen}{Group Knowledge Transfer in Neural Processes for Cold-Start Recommendation}
\newcommand{\authorname}{Владимир Лаврентьев}
\newcommand{\supervisorname}{Александр Геннадьевич Дьяконов}
\newcommand{\coursename}{Научная студия}
\newcommand{\placeholder}[1]{\textcolor{gray}{\textit{[#1]}}}

\hypersetup{
    pdftitle={\researchtitle},
    pdfauthor={\authorname}
}

\begin{document}
\pagenumbering{gobble}

% ---------- Титульный лист ----------
\begin{titlepage}
    \begin{center}
        \large Центральный университет

        \vspace{0.5cm}
        \large Курс «\coursename»

        \vfill

        \LARGE\textbf{\researchtitle}

        \vspace{0.5cm}
        \large \researchtitleen

        \vspace{1cm}
        \large Исследовательская работа
    \end{center}

    \vfill

    \begin{flushright}
        \begin{tabular}{@{}ll@{}}
            Автор: & \authorname \\
            Научный руководитель: & \supervisorname
        \end{tabular}
    \end{flushright}

    \vfill

    \begin{center}
        Москва, 2026
    \end{center}
\end{titlepage}
\pagenumbering{roman}

% ---------- Аннотации ----------
\section*{Аннотация}
\addcontentsline{toc}{section}{Аннотация}
% Ориентир: 200--250 слов, 5--7 предложений.
% Кратко отразите актуальность, цель, метод, основные результаты и вывод.
\placeholder{Текст аннотации на русском языке.}

\vspace{0.5cm}
\noindent\textbf{Ключевые слова:}
\placeholder{5--7 ключевых слов.}

\section*{Abstract}
\addcontentsline{toc}{section}{Abstract}
% Английская версия должна передавать тот же смысл, что и русская аннотация.
\placeholder{Abstract in English.}

\vspace{0.5cm}
\noindent\textbf{Keywords:}
\placeholder{5--7 keywords.}

\newpage
\tableofcontents
\newpage
\pagenumbering{arabic}

% ---------- Основной текст ----------
\section{Введение}
% Рекомендуемая логика введения:
% 1) мотивация и практическая значимость;
% 2) объект исследования и основные термины;
% 3) исследовательская проблема;
% 4) текущее состояние области;
% 5) цель и задачи;
% 6) предлагаемая идея и ожидаемая новизна;
% 7) ограничения существующих работ;
% 8) цель вычислительного эксперимента.
\placeholder{Почему задача важна, что именно исследуется и какой вклад планируется получить.}

\section{Обзор литературы}
% Сгруппируйте работы по подходам, а не пересказывайте статьи по одной.
% Покажите, что уже решено, какие ограничения остаются и где находится ваша работа.
\placeholder{Связанные исследования, базовые методы и незакрытый исследовательский вопрос.}

\section{Постановка задачи}
% Формально опишите:
% - данные и доступную информацию о пользователях и объектах;
% - режим cold-start и разбиение взаимодействий на support/query;
% - класс моделей или алгоритмов;
% - функцию обучения;
% - внешние метрики качества;
% - схему валидации и ограничения.
\placeholder{Формальная постановка задачи и обозначения.}

\section{Предлагаемый метод}
% Сначала дайте идею простыми словами, затем опишите алгоритм и его отличие от baseline.
% Для текущего проекта здесь можно отделить базовый TaNP от предлагаемого улучшения.
\placeholder{Описание базового метода, предлагаемого изменения и ожидаемого эффекта.}

\section{Вычислительный эксперимент}

\subsection{Данные и подготовка выборки}
% Укажите источник, размер, признаки, правила фильтрации и способ формирования задач.
\placeholder{Описание датасета и воспроизводимого препроцессинга.}

\subsection{Экспериментальный протокол}
% Зафиксируйте разбиение, baseline, конкурирующие методы, гиперпараметры,
% число запусков, random seeds, вычислительную среду и бюджет обучения.
\placeholder{Единые условия сравнения методов.}

\subsection{Метрики}
% Объясните, какую сторону качества измеряет каждая метрика и почему она нужна.
\placeholder{Основные и дополнительные метрики.}

\section{Результаты и обсуждение}
% Начните с главной таблицы сравнения. Затем добавьте анализ по числу взаимодействий,
% группам пользователей, абляции, устойчивости и вычислительной стоимости.
\placeholder{Результаты, их интерпретация, ограничения и возможные причины наблюдаемых эффектов.}

% Пример заготовки таблицы. Раскомментируйте, когда появятся результаты.
% \begin{table}[H]
%     \centering
%     \caption{Сравнение методов на тестовой выборке}
%     \label{tab:main-results}
%     \begin{tabular}{lcc}
%         \toprule
%         Метод & Метрика 1 & Метрика 2 \\
%         \midrule
%         Baseline & -- & -- \\
%         Предлагаемый метод & -- & -- \\
%         \bottomrule
%     \end{tabular}
% \end{table}

% Пример заготовки рисунка. Файлы рисунков храните в папке figures/.
% \begin{figure}[H]
%     \centering
%     \includegraphics[width=0.8\textwidth]{example.pdf}
%     \caption{Краткое содержательное название рисунка}
%     \label{fig:example}
% \end{figure}

\section{Заключение}
% Кратко ответьте на исследовательский вопрос, перечислите основные результаты,
% ограничения работы и следующие шаги. Не добавляйте здесь новые эксперименты.
\placeholder{Главные выводы и направления дальнейшей работы.}

% ---------- Библиография ----------
\clearpage
\addcontentsline{toc}{section}{Список литературы}
\bibliographystyle{plainnat}
\bibliography{references}

% ---------- Приложения ----------
\clearpage
\appendix
\section{Дополнительные материалы}
% Сюда можно вынести подробные гиперпараметры, дополнительные таблицы,
% схемы архитектуры и результаты вспомогательных экспериментов.
\placeholder{Материалы, которые важны для воспроизводимости, но перегружают основной текст.}

\end{document}
